Theorem 7. Let $J$ be a 1 -sphere in $\mathbf{R}^{2}$. Then $\mathbf{R}^{2}-J$ has only one bounded component.

Proof. Let $X, Y$, and $Z$ be as in Figure 4.2. Here, as usual, $J=A_{1} \cup A_{2} ; T$ is the first point of $B$ (in the order from $S$ ) that lies in $J, T \in A_{1} ; X$ is the last point of $B$ in $A_{1} ; Y$ is the first point after $X$ in $B$ that lies in $A_{2} ; Z$ is between $X$ and $Y$ in $B$; and $W$ is the last point of $B$ that lies in $A_{2}$. (See Figure 4.9.)

\begin{figure}[h]
\begin{center}
  \includegraphics[alt={},max width=\textwidth]{c570eb93-f29d-4347-b75d-c251ce2386f6-037_544_802_696_670}
\captionsetup{labelformat=empty}
\caption{Figure 4.9}
\end{center}
\end{figure}

We know that $Z$ lies in a bounded component of $\mathbf{R}^{2}-J$. We need to show that $\mathbf{R}^{2}-J$ has no other bounded component.

Let $B_{1}$ be the arc from $S$ to $T$ in $B$; let $B_{2}$ be the arc from $T$ to $X$ in $A_{1}$ (if $T$ and $X$ are really different; if not, let $B_{2}=T=X$ ); let $B_{3}$ be the arc from $X$ to $Y$ in $B$; let $B_{4}$ be the arc from $Y$ to $W$ in $A_{2}$ (if $Y \neq W$ ); and let $B_{5}$ be the arc from $W$ to $Q$ in $B$. Let $B^{\prime}=\bigcup_{i=1}^{5} B_{i}$. Then $P$ and $R$ are limit points of different components of $I-B^{\prime}$. Therefore, if $U$ is a bounded component of $\mathbf{R}^{2}-J$, different from the component which contains $Z$, it follows that $U \cap B^{\prime}=\varnothing$, so that $\operatorname{Fr} U$ cannot contain both $P$ and $R$. Therefore $\mathrm{Fr} U$ lies in an $\operatorname{arc}$ in $J$; and this is impossible, by Theorem 5. \(\square\)
